\section{Metrik dan Instrumen Evaluasi}

Definisi konseptual variabel terikat telah dijelaskan sebelumnya. 
Bagian ini menetapkan definisi operasional, formula, dan instrumen pengukuran setiap metrik yang digunakan dalam eksperimen.

Akurasi kebijakan diukur menggunakan \textit{F1-score}, yaitu rata-rata harmonik antara presisi dan \textit{recall} keputusan akses yang dihasilkan kebijakan hasil \textit{mining} terhadap \textit{log} akses uji atau \textit{holdout}. 
F1\label{first:f1} didefinisikan sebagai berikut:

\[
F_1 = 2 \times \frac{\mathrm{Precision} \times \mathrm{Recall}}{\mathrm{Precision} + \mathrm{Recall}},
\]
\[
\quad
\text{dengan }
\mathrm{Precision} = \frac{\mathrm{TP}}{\mathrm{TP} + \mathrm{FP}}
\quad
\text{dan }
\mathrm{Recall} = \frac{\mathrm{TP}}{\mathrm{TP} + \mathrm{FN}}.
\]
\label{first:precision}\label{first:recall}

Dengan \textit{permit} sebagai kelas positif, TP\label{first:tp} (\textit{true positive}) adalah permintaan yang seharusnya diizinkan dan diprediksi \textit{permit}, FP\label{first:fp} (\textit{false positive}) adalah permintaan yang seharusnya ditolak tetapi diprediksi \textit{permit}, sedangkan FN\label{first:fn} (\textit{false negative}) adalah permintaan yang seharusnya diizinkan tetapi diprediksi \textit{deny}.

Kompleksitas struktural kebijakan diukur menggunakan \textit{Weighted Structural Complexity} (WSC), yaitu jumlah pasangan atribut–nilai yang digunakan pada seluruh aturan kebijakan. 
WSC dirumuskan sebagai berikut:

\[
\mathrm{WSC}(\pi) = \sum_{r \in \pi} \lvert \mathrm{attr}(r) \rvert,
\]\label{first:symbol-r}

dengan $\pi$ menyatakan kebijakan yang dihasilkan dan $|attr(r)|$\label{first:symbol-attrr} jumlah pasangan atribut–nilai pada aturan $r$. 
Secara \textit{default}, seluruh pasangan atribut–nilai diberi bobot seragam sebesar 1, meskipun pembobotan berbeda untuk atribut subjek, objek, atau lingkungan dapat diterapkan bila diperlukan.

\textit{Coverage} mengukur proporsi permintaan akses pada \textit{log} uji yang dicakup oleh sedikitnya satu aturan dalam kebijakan hasil \textit{mining}, sehingga tidak bergantung pada keputusan \textit{default} seperti \textit{default-deny}. 
Metrik ini digunakan untuk menilai kelengkapan cakupan kebijakan:

\[
\mathrm{Coverage} = \frac{\text{jumlah permintaan akses yang cocok dengan }\label{first:coverage} \ge 1 \text{ aturan}}{\text{total permintaan akses pada log uji}}.
\]

F1, WSC, dan \textit{coverage} secara bersama merepresentasikan kualitas kebijakan yang digunakan pada Eksperimen 1.

Jejak memori diukur sebagai konsumsi memori puncak selama eksekusi dan dilaporkan dalam dua bentuk. 
$MemPeak_{alloc}$ adalah puncak alokasi struktur data algoritma yang diukur menggunakan $tracemalloc$ dan digunakan sebagai metrik utama, sedangkan $MemPeak_{sys}$ adalah nilai VmHWM pada tingkat sistem operasi yang merepresentasikan puncak RSS\label{first:rss} proses. 
Selain nilai absolut, memori dianalisis terhadap ukuran populasi virtual (NP) dan dimensi (D). 
Bentuk kurva penskalaan ini digunakan pada Eksperimen 2 sebagai bukti empiris pendukung H2, yaitu bahwa kebutuhan memori CoDA-EDA tidak sensitif terhadap $NP$ dan bertumbuh secara linear terhadap $D$.

Waktu komputasi dibedakan menjadi $T_{init}$\label{first:tinit}, yaitu waktu pembelajaran awal; $T_{param}$\label{first:tparam}, yaitu waktu siklus yang hanya memperbarui parameter \textit{Probability Vector}; dan $T_{struct}$\label{first:tstruct}, yaitu waktu siklus yang juga melakukan pembelajaran ulang struktur setelah \textit{drift} terdeteksi. 
Ketiganya diukur menggunakan \textit{wall-clock time} pada perangkat target dan digunakan pada Eksperimen 3 untuk membandingkan efisiensi pembaruan CoDA-EDA dan iBOA.

Pada skenario \textit{drift}, efektivitas strategi \textit{warm-start} diukur melalui dua metrik tambahan. 
$T_{reopt}$\label{first:treopt} menyatakan waktu sejak \textit{drift} terdeteksi hingga kebijakan memenuhi kriteria konvergensi, sedangkan $\Delta Quality$\label{first:delta-quality} menyatakan perbedaan F1 dan WSC antara \textit{warm-start} dan \textit{re-mining} penuh pada titik evaluasi yang setara. 
Kedua metrik tersebut digunakan untuk menilai apakah \textit{warm-start} dapat mempertahankan kualitas kebijakan dengan waktu re-optimasi yang lebih rendah.

\begin{longtable}{|c|L{2.3cm}|L{2.6cm}|L{3.3cm}|L{3.3cm}|p{1.6cm}|}
\caption{Definisi Operasional Metrik dan Instrumen Evaluasi}
\label{tab:definisi-operasional-metrik} \\
\hline
\textbf{No} & \textbf{Metrik} & \textbf{Simbol} & \textbf{Definisi/Formula} & \textbf{Instrumen} & \textbf{Hipotesis} \\
\hline
\endfirsthead

\multicolumn{6}{c}{{\bfseries \tablename\ \thetable{} -- lanjutan dari halaman sebelumnya}} \\
\hline
\textbf{No} & \textbf{Metrik} & \textbf{Simbol} & \textbf{Definisi/Formula} & \textbf{Instrumen} & \textbf{Hipotesis} \\
\hline
\endhead

\hline
\multicolumn{6}{r}{Lanjut ke halaman berikutnya} \\
\endfoot

\hline
\endlastfoot

% ==================== DATA BARIS ====================
1 & Akurasi & F1 & Rata-rata harmonik presisi dan recall keputusan akses terhadap log uji & Skrip evaluasi Python terhadap log uji holdout & H1 \\
\hline
2 & Kompleksitas struktural & WSC & Jumlah pasangan atribut-nilai pada seluruh aturan kebijakan & Skrip evaluasi Python & H1 \\
\hline
3 & Cakupan & Coverage & Proporsi permintaan akses yang tercakup $\ge$ 1 aturan & Skrip evaluasi Python & H1 \\
\hline
4 & Jejak memori & MemPeak\_alloc (utama); MemPeak\_sys (konteks) & Puncak alokasi struktur data algoritma; dan puncak resident set size pada tataran sistem. Dilaporkan pula sebagai kurva penskalaan terhadap NP dan D & tracemalloc (utama) + VmHWM pada /proc/[pid]/status, satu proses terpisah per konfigurasi (Subbab 3.7) & H2 \\
\hline
5 & Waktu inisialisasi & T\_init & Waktu hingga elite pertama terbentuk & Wall-clock timer pada perangkat target & H3 \\
\hline
6 & Waktu siklus parameter & T\_param & Waktu satu siklus pembaruan Probability Vector saja & Wall-clock timer pada perangkat target & H3 \\
\hline
7 & Waktu siklus struktur & T\_struct & Waktu satu siklus dengan pembaruan struktur dependensi & Wall-clock timer pada perangkat target & H3 \\
\hline
8 & Waktu re-optimasi & T\_reopt & Waktu sejak drift terdeteksi hingga konvergensi tercapai & Wall-clock timer pada perangkat target & H4 \\
\hline
9 & Selisih kualitas warm-start & $\Delta$Quality & Selisih F1/WSC antara warm-start dan re-mining penuh & Skrip evaluasi Python & H4 \\
\hline

\end{longtable}

Definisi operasional dan formula pada Tabel \ref{tab:definisi-operasional-metrik} menjadi landasan pengolahan data dari pelaksanaan demonstrasi, sekaligus menjadi rujukan bagi teknik analisis statistik yang akan dilakukan.